\topic{先分清地址范围和已经驻留的页面}

执行 \code{exec("/bin/tool")} 时，内核先读取 ELF 程序头。程序头描述代码、
只读数据和可写数据应出现在哪些虚拟地址，却没有要求内核立刻为每一页分配
物理内存。项目用 VMA 保存这份意图，用页表保存当前已经能够访问的页面。

\begin{osgridlongtable}{L{0.22\textwidth}L{0.32\textwidth}L{0.36\textwidth}}
\hline
\textbf{结构} & \textbf{回答的问题} & \textbf{例子}\\ \hline
VMA & 这段地址应该存在吗，允许怎样访问，数据从哪里来 &
\code{0x400000--0x404000} 是只读可执行的 ELF 映像\\ \hline
PTE & 这一页此刻映射到哪个物理页帧，硬件允许怎样访问 &
\code{0x401000} 已映射为 user/read/execute\\ \hline
页帧引用 & 多个地址空间是否还共享同一物理页 & fork 后父子页表共同引用
原数据页\\ \hline
\end{osgridlongtable}

用户第一次访问尚未驻留的地址时，CPU 产生页故障并把地址写入 CR2。内核先用
VMA 判断访问是否合理。匿名区域得到一页清零内存；文件后备区域从对应文件
offset 读取内容；落在区域之外或违反权限的访问终止当前用户进程。

\topic{exec 为什么要先建立候选映像}
装入新程序会改动页表、VMA、用户栈、参数和入口寄存器。若在中途直接覆盖当前
进程，第五次分配失败时，旧程序已经被拆掉，新程序又不能运行。项目先在独立
地址空间中准备这些对象，全部检查通过后才替换当前映像。失败路径销毁候选
地址空间，调用者仍能继续执行旧程序。

\topic{fork 为什么先共享页面}
逐页复制整个地址空间会让 \code{fork} 的时间和内存消耗随进程大小增长，而
子进程常常紧接着调用 \code{exec}，许多复制没有被使用。项目把父子页表中的
可写页暂时改成只读，并增加物理页引用。某一方第一次写入时触发页故障；内核
分配新页、复制旧内容，只修改发起写入的地址空间。

可观察结果比“fork 成功”更具体。运行
\filepath{source/user/programs/fork\_probe.cpp} 后，应看到
\code{COW_ISOLATION_VERIFIED}：子进程写入私有页没有改变父进程看到的字节。
\filepath{source/user/programs/memory\_probe.cpp} 另外检查按需清零、栈增长、
文件映射和共享写回。后面的四节会沿这些现象分别进入进程树、匿名 VMA、
文件后备 VMA 与 COW 实现。
